\chapter{绪论}\label{Chap:chap1}


\section{研究背景与意义}\label{Sec1.1}

强化学习（Reinforcement Learning, RL）是一种实现机器自主决策的人工智能范式，智能体通过与环境的持续交互学习最优策略，以最大化长期累积回报\cite{Levine2020OfflineRL}。自深度强化学习在围棋\cite{Silver2016AlphaGo}等领域取得突破以来，学术界与工业界持续将其应用拓展至更复杂的真实场景。近年来，强化学习在机器人控制\cite{Johannink2019Residual}领域积累了丰富经验，并随着DeepSeek-R1\cite{DeepSeekAI2025R1}在推理任务中的成功应用，开始广泛用于大语言模型的对齐与微调。

在线强化学习框架中，智能体直接与环境交互，通过试错优化行为策略。该方法具备良好的适应性，能够根据环境反馈及时调整策略，但需要大量交互样本。在医疗诊断\cite{Gottesman2019Healthcare}、自动驾驶\cite{Kiran2022Driving}等数据采集成本高昂或存在安全约束的场景中，频繁的在线试错既不经济也存在风险，限制了在线强化学习的实际部署。

离线强化学习提供了一种替代方案：从预先收集的静态数据集中学习策略，避免在线交互的成本与风险。该范式依赖行为策略采集的历史数据，训练开销较小。然而，行为策略通常是次优的，且存在分布偏移问题\cite{Fujimoto2019BCQ}——训练数据分布与目标策略访问的状态-动作分布存在差异，导致价值函数对分布外动作的估计不准确。由于离线学习无法获得环境反馈修正估计误差，累积误差会引发对分布外动作价值的过高估计，产生外推误差\cite{Liu2020Provably}。如何从有限且可能存在偏差的数据中学习泛化能力良好的策略，是离线强化学习的核心挑战。

离线到在线强化学习结合了两种范式的优势：先在离线数据上预训练策略，再通过有限的在线交互进行微调，既利用了历史数据降低样本需求，又通过在线反馈提升策略性能。

当前离线到在线方法多采用无模型框架，侧重于离线与在线数据的融合策略，如平衡重放、优先级采样等。无模型方法的样本效率受限，难以充分挖掘大规模离线数据集的价值。本文研究基于不确定性探索模型的离线到在线强化学习算法，将优化重点置于环境动力学模型：利用离线策略与环境交互获取的在线数据优化状态转移模型，构建更准确且覆盖分布外区域的环境模型，进而完成策略更新。该方法以模型生成数据替代频繁的环境交互，在线开销较小，数据利用效率较高。

引入不确定性量化机制后，该方法在离线阶段可识别数据集中的高不确定性区域，对这些区域采取保守的策略更新；在线阶段则优先探索不确定性较高的状态-动作对，以较小的交互代价获取策略未覆盖区域的新数据。这种不确定性引导的探索机制提升了模型对环境变化的适应能力。

该算法适用于需要充分利用历史数据、同时在线探索受限的场景，如自动驾驶、金融决策、医疗优化等。基于模型的框架通过数据生成减少对在线交互的依赖，不确定性引导的探索策略提升了在线微调的样本效率，在控制在线成本的同时实现策略性能的持续提升。


\section{国内外研究现状}\label{Sec1.2}

强化学习在决策优化领域的应用日趋成熟，学术界与工业界均取得了阶段性进展。在大语言模型领域，DeepSeek发布的R1系列模型\cite{DeepSeekAI2025R1}通过纯强化学习范式激发模型的复杂推理能力，验证了强化学习在语言模型对齐中的有效性；ByteDance等团队提出的DAPO算法\cite{Yu2025DAPO}开源了大规模语言模型强化学习训练系统，在Qwen2.5-32B基座模型上实现了可复现的推理能力训练。在具身智能领域，Google DeepMind研发的RT-2模型\cite{Brohan2023RT2}将视觉-语言预训练知识迁移至机器人控制，提升了机器人在真实环境中的语义理解与操作泛化能力。在离线强化学习理论研究方面，王雪松等\cite{Wang2024Representation}系统梳理了基于表征学习的离线强化学习方法，分析了特征解耦对策略迁移能力的增强作用。

上述进展表明强化学习在大模型对齐、表征学习、复杂场景适配等方向取得了实质性成果，但也对离线到在线强化学习技术提出了更高要求。当前该领域仍面临若干技术挑战：离线阶段基于固定数据集训练的策略在在线环境中因数据分布差异易产生性能衰减，而在线微调过程中又可能丢失离线阶段习得的知识，现有方案通过动态策略调整、重要性加权等机制尝试平衡两者，但需针对具体任务分析数据分布特性并设计回放机制，开发成本较高。此外，探索-利用权衡是强化学习的固有难题，离线到在线的两阶段模式进一步放大了这一矛盾——离线阶段需保持策略保守性以规避数据偏差风险，在线阶段又需探索未知区域以适配环境变化，两阶段对探索强度的需求相反，RT-2模型虽提升了泛化能力，但未形成适配两阶段特性的动态探索机制。在数据利用效率方面，离线数据集的知识常未被充分利用，策略在微调阶段易丢失预训练知识；在线交互虽能获取环境反馈，但频繁访问环境会削弱离线学习低成本的优势，如何在提升数据利用率的同时控制在线开销，是该领域亟待解决的问题。

上述挑战推动了基于不确定性探索等技术方案的研究，也是本文的研究出发点。


\section{研究内容与主要工作}\label{Sec1.3}
% TODO: 根据论文实际内容补充


\section{文章结构安排}\label{Sec1.4}
% TODO: 根据论文实际章节安排补充






























